\documentclass[a4paper]{article}
\usepackage{martimath}
\title{Estructures Algebraiques}
\author{Martí Carrasco}
\date{}

\begin{document}
\maketitle
\tableofcontents
\indexresultats
\clearpage

\section{Anells i ideals}
\subsection{Anell quocient}

\begin{definicio}[Ideal]
Sigui $A$ un anell commutatiu. Un subconjunt $I\subseteq A$
és un ideal si és un subgrup additiu i $ax\in I$ per a tot
$a\in A$ i $x\in I$.
\end{definicio}

\begin{definicio}[Anell quocient][separat]
Sigui $I$ un ideal de $A$. El conjunt $A/I$ de classes
d'equivalència té una estructura natural d'anell definida per
\[
\overline a+\overline b=\overline{a+b},\qquad
\overline a\,\overline b=\overline{ab}.
\]
\end{definicio}

\begin{obs}[][separat]
La projecció canònica $\pi:A\to A/I$ és un morfisme
exhaustiu amb nucli $I$.
\end{obs}

\begin{teo}[Primer teorema d'isomorfisme][important,separat]
Sigui $f:A\to B$ un morfisme d'anells. Aleshores,
\[
A/\ker f\cong\operatorname{Im} f.
\]
\begin{dem}
Definim $\overline f(\overline a)=f(a)$.
Si $\overline a=\overline b$, aleshores $a-b\in\ker f$,
i per tant $f(a)=f(b)$. L'aplicació induïda és bijectiva.
\end{dem}
\end{teo}

\propi[Ideals principals]{Si $a\in A$, el conjunt
$(a)=\{ra:r\in A\}$ és un ideal de $A$.}

\begin{exemple}[Els enters mòdul $6$][separat]
En $\mathbb Z/6\mathbb Z$ tenim
\[
\overline 2\,\overline 4=\overline 2,\qquad
\overline 3+\overline 5=\overline 2.
\]
\end{exemple}

\subsubsection{Una propietat útil}

\begin{lema}[Nucli additiu]
El nucli d'un morfisme d'anells és un subgrup additiu.
\end{lema}
\begin{demostracio}
Com que tot morfisme d'anells és additiu, el seu nucli
és un subgrup del grup additiu corresponent.
\end{demostracio}

\begin{cor}[Conseqüència]
El nucli d'un morfisme d'anells commutatius és un ideal.
\end{cor}

\begin{comentari}
Els títols són opcionals; per defecte el text és en línia.
\end{comentari}

\section{Figures i gràfiques}
\begin{grafica}[xmin=-2.2,xmax=2.2,ymin=-1.5,ymax=3.4]
\plotfuncion[color=mmblue]{x^2-1}
\plotfuncion[color=black,dashed]{x+0.2}
\end{grafica}
\begin{center}\footnotesize Gràfica generada amb \texttt{pgfplots}.\end{center}

\end{document}
